\chapter{Transfer-Operator Gaps and Block Separation}\label{ch:spectral}

This chapter proves that the MPV overlaps between distinct
(non-gauge-phase-equivalent) MPS blocks decay to zero as system
size grows. It also proves the corresponding self-overlap convergence
to~$1$ under the complementary transfer-map gap hypothesis that encodes
primitivity in the cases of interest. The argument is based on a
spectral analysis of the \emph{mixed transfer operator},
following~\cite{PerezGarcia2007Matrix}.

This chapter concerns gaps in the spectrum of transfer operators,
writing $\rho_{\spec}(\cdot)$ for the spectral radius:
$\rho_{\spec}(F_{AB}) < 1$ for mixed transfer operators, and
$\rho_{\spec}(\E_A - P) < 1$ after removing the fixed-point projection. The
unqualified phrase \emph{spectral gap} is reserved for the energy gap of a
many-body Hamiltonian in Chapter~\ref{ch:parent_hamiltonian}.

The algebraic mixed-transfer identities in the first two sections do
not use normalization. Starting with the spectral-radius estimates,
we assume the trace-preserving normalization
\begin{equation}\label{eq:spec_tp_normalization}
    \sum_i (A^i)^\dagger A^i = \Id.
\end{equation}

\section{Mixed transfer operator}

\begin{definition}[Mixed transfer operator]\label{def:mixed_transfer}
    \lean{MPSTensor.mixedTransferMap}
    \leanok
    \uses{def:mps_tensor}
    The \emph{mixed transfer operator} (or \emph{cross transfer
        operator}) for two MPS tensors $A$ and $B$ of the same bond
    dimension~$D$ is the linear map
    $F_{AB} : \MN{D} \to \MN{D}$ defined by
    \[
        F_{AB}(X) =  \sum_{i=0}^{d-1} A^i X (B^i)^\dagger.
    \]
\end{definition}

\begin{definition}[Rectangular mixed transfer operator]\label{def:mixed_transfer_rect}
    \lean{MPSTensor.mixedTransferMap₂}
    \leanok
    \uses{def:mps_tensor}
    For MPS tensors $A$ of bond dimension $D_1$ and $B$ of bond
    dimension $D_2$ (possibly $D_1 \neq D_2$), the
    \emph{rectangular mixed transfer operator} is
    $F_{AB}^{\mathrm{rect}} :
        M_{D_1 \times D_2}(\C) \to M_{D_1 \times D_2}(\C)$
    defined by
    \[
        F_{AB}^{\mathrm{rect}}(X)
        =  \sum_{i=0}^{d-1} A^i X (B^i)^\dagger.
    \]
\end{definition}

\begin{theorem}[Self-transfer]\label{thm:mixed_self}
    \lean{MPSTensor.mixedTransferMap_self}
    \leanok
    \uses{def:mixed_transfer, def:transfer_map}
    Setting $B=A$ recovers the ordinary transfer map: $F_{AA} = \E_A$.
\end{theorem}
\begin{proof}\leanok
    Substituting $B=A$ into the defining sum
    $F_{AB}(X) = \sum_i A^i X (B^i)^\dagger$ gives
    $\sum_i A^i X (A^i)^\dagger = \E_A(X)$.
\end{proof}

\begin{lemma}[Iterated mixed transfer]\label{thm:mixed_pow}
    \lean{MPSTensor.mixedTransferMap_pow_apply}
    \leanok
    \uses{def:mixed_transfer, def:eval_word}
    For all $X \in \MN{D}$ and $N \ge 0$,
    \[
        F_{AB}^N(X)
        =  \sum_{\sigma \in \{0,\ldots,d{-}1\}^N}
        A^\sigma X (B^\sigma)^\dagger,
    \]
    where $A^\sigma = A^{\sigma_1} \cdots A^{\sigma_N}$ denotes the
    word evaluation (Definition~\ref{def:eval_word}).
\end{lemma}

\begin{proof}
    \leanok
    Induction on $N$: the base case is $F^0(X) = X$;
    the inductive step applies $F_{AB}$ to each summand
    and re-indexes via
    $\{0,\ldots,d{-}1\}^{N+1} \cong
        \{0,\ldots,d{-}1\} \times \{0,\ldots,d{-}1\}^N$.
\end{proof}

\begin{lemma}[Matrix trace of iterated mixed transfer at the identity]\label{thm:mixed_trace_identity}
    \lean{MPSTensor.trace_mixedTransferMap_pow_identity}
    \leanok
    \uses{def:mixed_transfer, def:eval_word}
    For $N \ge 0$,
    \[
        \tr(F_{AB}^N(\Id))
        =
        \sum_{\sigma}
        \tr(A^\sigma (B^\sigma)^\dagger).
    \]
    This is the matrix trace of $F_{AB}^N$ evaluated at the
    identity matrix.
    It should \emph{not} be confused with the operator trace
    $\Tr(F_{AB}^N)$; its relation to the MPV overlap is proved in the
    next section.
\end{lemma}

\begin{proof}\leanok\uses{thm:mixed_pow}
    Apply the matrix trace to each summand in
    Lemma~\ref{thm:mixed_pow} with $X = \Id$:
    $\tr(F_{AB}^N(\Id))
        = \sum_\sigma \tr(A^\sigma \cdot \Id \cdot (B^\sigma)^\dagger)
        = \sum_\sigma \tr(A^\sigma (B^\sigma)^\dagger)$.
\end{proof}

\section{MPV overlap as transfer trace}

Recall from Definition~\ref{def:mpv_overlap} that
\[
    O_{AB}(N) = \sum_{\sigma \in \{0,\ldots,d{-}1\}^N}
    V^{(N)}(A)_\sigma \overline{V^{(N)}(B)_\sigma}.
\]
This section rewrites that scalar as the operator trace of the
mixed transfer power.

\begin{lemma}[Trace expansion over matrix units]\label{lem:trace_expansion}
    \lean{MPSTensor.linearMap_trace_eq_sum_apply_single}
    \leanok
    For any linear endomorphism $T$ on $M_{D_1 \times D_2}(\C)$,
    the operator trace can be expanded as
    \[
        \Tr(T)
        =
        \sum_{p=0}^{D_1-1} \sum_{q=0}^{D_2-1}
        (T(E_{pq}))_{pq},
    \]
    where $E_{pq} \in M_{D_1 \times D_2}(\C)$ is the matrix with
    $1$ in position $(p,q)$ and $0$ elsewhere.
\end{lemma}

\begin{proof}\leanok
    The operator trace $\Tr(T)$ equals the sum of
    $\tr(T(E_{pq}) \cdot E_{qp})$ over all $(p,q)$,
    and $\tr(M \cdot E_{qp}) = M_{pq}$.
\end{proof}

\begin{lemma}[Overlap equals transfer trace]\label{thm:overlap_trace}
    \lean{MPSTensor.trace_mixedTransferMap_pow_eq_mpvOverlap}
    \leanok
    \uses{def:mixed_transfer, def:mpv_overlap}
    For tensors $A, B$ of bond dimension~$D \ge 1$,
    \[
        O_{AB}(N)
        =
        \Tr(F_{AB}^N).
    \]
\end{lemma}

\begin{proof}
    \leanok
    \uses{thm:mixed_pow, lem:trace_expansion}
    Expand the operator trace using
    Lemma~\ref{lem:trace_expansion} as a sum over matrix units
    $E_{pq}$.
    Apply Lemma~\ref{thm:mixed_pow} to evaluate
    $F_{AB}^N(E_{pq})$, extract the $(p,q)$-entry, and
    reassemble. The resulting double sum over $(p,q)$ and $\sigma$
    equals $\sum_\sigma V^{(N)}(A)_\sigma
        \overline{V^{(N)}(B)_\sigma} = O_{AB}(N)$.
\end{proof}

\begin{lemma}[Rectangular overlap as transfer trace]\label{thm:overlap_trace_rect}
    \lean{MPSTensor.trace_mixedTransferMap₂_pow_eq_mpvOverlap}
    \leanok
    \uses{def:mixed_transfer_rect, def:mpv_overlap}
    For tensors $A$ of bond dimension $D_1 \ge 1$ and $B$ of bond
    dimension $D_2 \ge 1$,
    \[
        O_{AB}(N) = \Tr((F_{AB}^{\mathrm{rect}})^N).
    \]
\end{lemma}

\begin{proof}\leanok
    \uses{thm:overlap_trace}
    Expand the operator trace over the rectangular matrix-unit basis.
    Then apply the rectangular iterated-transfer formula and sum over
    the matrix-unit indices exactly as in
    Lemma~\ref{thm:overlap_trace}.
\end{proof}

\begin{lemma}[Word trace as entrywise inner product]\label{lem:word_trace_entrywise}
    \lean{MPSTensor.mpv_inner_product_via_trace}
    \leanok
    \uses{def:eval_word}
    For a word $\sigma \in \{0,\ldots,d{-}1\}^N$,
    \[
        \tr(A^\sigma (B^\sigma)^\dagger)
        =
        \sum_{j,k = 0}^{D-1}
        (A^\sigma)_{jk} \overline{(B^\sigma)_{jk}}.
    \]
\end{lemma}

\begin{proof}\leanok
    Expand the matrix trace and the $(j,j)$-entry of
    $A^\sigma (B^\sigma)^\dagger$.
\end{proof}

\section{Frobenius norm estimates}

The transfer-operator gap proofs use the Frobenius (Hilbert--Schmidt) norm
as the natural norm for the finite-dimensional matrix algebra.
We write $\|{\cdot}\|_F$ for this norm; it coincides with, and we use it
interchangeably with, the Hilbert--Schmidt norm $\|{\cdot}\|_{\mathrm{HS}}$ on finite-dimensional spaces.
We use the squared version and an isometric embedding into
Euclidean space.

\begin{definition}[Frobenius norm squared]\label{def:frob_sq}
    \lean{MPSTensor.frobSq}
    \leanok
    For a (possibly rectangular) matrix $X \in M_{m \times n}(\C)$,
    the \emph{Frobenius norm squared} is
    \[
        \|X\|_F^2 := \sum_{i=0}^{m-1}\sum_{j=0}^{n-1} |X_{ij}|^2.
    \]
\end{definition}

\begin{lemma}[Trace formula for Frobenius norm]\label{lem:frob_sq_trace}
    \lean{MPSTensor.frobSq_trace}
    \leanok
    \uses{def:frob_sq}
    $\|X\|_F^2 = \operatorname{Re}\tr(X^\dagger X)$.
\end{lemma}

\begin{proof}\leanok
    Expand the matrix product and trace entry by entry.
\end{proof}

\begin{definition}[Euclidean-space embedding]\label{def:mat_to_es}
    \lean{MPSTensor.matToES}
    \leanok
    The map $\mathrm{vec} : M_{m \times n}(\C) \to \C^{mn}$ that
    flattens matrix entries is an isometric embedding with respect
    to the Frobenius norm: $\|\mathrm{vec}(X)\|^2 = \|X\|_F^2$.
\end{definition}

\begin{lemma}[Norm of embedded matrix]\label{lem:norm_mat_to_es}
    \lean{MPSTensor.norm_matToES_sq}
    \leanok
    \uses{def:mat_to_es, def:frob_sq}
    $\|\mathrm{vec}(X)\|^2 = \|X\|_F^2$.
\end{lemma}

\begin{proof}\leanok
    Expand both sides as sums over entries and use $|z|^2 = \bar{z}z$.
\end{proof}

\section{Eigenvalue bound and transfer-operator gap}

The eigenvalue bound comes from a uniform Frobenius-norm estimate
obtained directly from the trace-preserving normalization
(cf.~\cite[Proposition~6.1]{Wolf2012Quantum} for the general
operator-norm version via Russo--Dye).
For the rigidity theorem, one gauges $A$ and $B$ separately using
positive fixed points of their transfer maps to obtain unital Kraus
families, and then applies the Kadison--Schwarz equality argument
(\cite[Theorem~5.3]{Wolf2012Quantum}) to
a block embedding of a mixed-transfer eigenvector.

\begin{theorem}[Eigenvalue bound]\label{thm:eigenvalue_bound}
    \lean{MPSTensor.eigenvalue_norm_le_one}
    \leanok
    \uses{def:mixed_transfer}
    Let $A, B$ be normalized MPS tensors.
    Every eigenvalue $\mu$ of $F_{AB}$ satisfies $|\mu| \le 1$.
\end{theorem}

\begin{proof}
    \leanok
    \uses{thm:mixed_pow}
    Expand $F_{AB}^n(X)$ as a sum over words and use Cauchy--Schwarz,
    together with
    $\sum_i (A^i)^\dagger A^i = \sum_i (B^i)^\dagger B^i = \Id$
    (Equation~(\ref{eq:spec_tp_normalization}), applied to both tensors),
    to obtain a uniform Frobenius-norm bound
    $\|F_{AB}^n(X)\|_{\mathrm{HS}} \le C_D \|X\|_{\mathrm{HS}}$
    for all $n$, where $C_D$ depends only on $D$.
    If $F_{AB}(X) = \mu X$ and $X \neq 0$, then
    $|\mu|^n \|X\|_{\mathrm{HS}}
        = \|F_{AB}^n(X)\|_{\mathrm{HS}}
        \le C_D \|X\|_{\mathrm{HS}}$
    for every $n$.
    Hence $|\mu| \le 1$.

    The argument uses only the trace-preserving normalization
    (\ref{eq:spec_tp_normalization})
    and does not appeal to Perron--Frobenius theory.
\end{proof}

\begin{theorem}[Rectangular eigenvalue bound]\label{thm:eigenvalue_bound_rect}
    \lean{MPSTensor.eigenvalue_norm_le_one₂}
    \leanok
    \uses{def:mixed_transfer_rect}
    Let $A$ and $B$ be normalized MPS tensors of bond dimensions
    $D_1 \ge 1$ and $D_2 \ge 1$.  Every eigenvalue $\mu$ of
    $F_{AB}^{\mathrm{rect}}$ satisfies $|\mu| \le 1$.
\end{theorem}

\begin{proof}\leanok
    The word-expansion and Frobenius-norm estimate used in
    Theorem~\ref{thm:eigenvalue_bound} apply equally to rectangular matrices:
    if $F_{AB}^{\mathrm{rect}}(X)=\mu X$ and $X\neq 0$, then
    $|\mu|^n\|X\|_{\mathrm{HS}}\le D_1\|X\|_{\mathrm{HS}}$
    for every $n$.  Hence $|\mu|\le 1$.
\end{proof}

\begin{lemma}[Spectral radius bound]\label{thm:spectral_radius_le_one}
    \lean{MPSTensor.spectralRadius_mixedTransfer_le_one}
    \leanok
    \uses{def:mixed_transfer}
    $\rho_{\spec}(F_{AB}) \le 1$ for normalized tensors.
\end{lemma}

\begin{proof}\leanok\uses{thm:eigenvalue_bound}
    The spectral radius is the supremum of $|\mu|$ over all
    eigenvalues, and each satisfies $|\mu| \le 1$ by
    Theorem~\ref{thm:eigenvalue_bound}.
\end{proof}

\begin{theorem}[Spectral radius $\ge 1$ implies gauge-phase equivalence]%
    \label{thm:modulus_one_gauge}
    \lean{MPSTensor.modulus_one_eigenvalue_implies_gauge}
    \leanok
    \uses{def:mixed_transfer, def:gauge_phase_equiv, def:injective}
    Let $A, B$ be injective normalized MPS tensors.
    If $\rho_{\spec}(F_{AB}) \ge 1$, then $A$ and $B$ are
    gauge-phase equivalent.
\end{theorem}

\begin{proof}
    \leanok
    \uses{thm:spectral_radius_le_one, def:gauge_phase_equiv,
        thm:qpf, thm:right_canonical_gauge,
        thm:ks_peripheral, thm:kraus_commute,
        thm:left_multiplicative_identity, thm:simplicity, thm:skolem_noether}
    Since $\rho_{\spec}(F_{AB}) \le 1$ (Lemma~\ref{thm:spectral_radius_le_one})
    and $\rho_{\spec}(F_{AB}) \ge 1$ by hypothesis, there exists an
    eigenvalue~$\mu$ with $|\mu| = 1$ and eigenvector $X \neq 0$:
    $F_{AB}(X) = \mu X$.
    The proof proceeds in five steps.

    \emph{Step~1 (Separate gauging).}
    Choose positive definite fixed points $\rho_A, \rho_B$ of the
    transfer maps $\E_A, \E_B$ (Theorem~\ref{thm:qpf}) and gauge each
    tensor separately to a unital Kraus family
    (Theorem~\ref{thm:right_canonical_gauge}): set
    $A'^i = \rho_A^{-1/2} A^i \rho_A^{1/2}$ and similarly for~$B$.
    Let $X' = \rho_A^{-1/2} X \rho_B^{-1/2}$ denote the correspondingly
    gauged eigenvector.

    \emph{Step~2 (Block Kraus family and off-diagonal embedding).}
    Form the block Kraus operators $C^i = \bigl(\begin{smallmatrix}
            A'^i & 0 \\ 0 & B'^i\end{smallmatrix}\bigr)$.
    These constitute a unital Kraus family on $\MN{2D}$.
    The off-diagonal matrix
    $Y = (\begin{smallmatrix} 0 & X' \\ 0 & 0\end{smallmatrix})$
    satisfies $E_C(Y) = \mu Y$.

    \emph{Step~3 (KS equality and Kraus commutation).}
    Because $|\mu| = 1$ and $E_C$ is unital and trace-preserving,
    the Kadison--Schwarz equality for peripheral eigenvectors
    (Theorem~\ref{thm:ks_peripheral}) gives
    $E_C(Y^\dagger Y) = E_C(Y)^\dagger E_C(Y)$.
    By the Kraus commutation relation
    (Theorem~\ref{thm:kraus_commute}), $Y$ commutes with every block
    Kraus operator: $C^i Y = Y C^i$ for all~$i$.

    \emph{Step~4 (Intertwining identity and invertibility).}
    Expanding the commutation relation block by block gives
    $A'^i X' = X' B'^i$ for every~$i$.
    Since $\{A'^i\}$ span $\MN{D}$ (injectivity of~$A$), $X'$ is a
    nonzero intertwiner from $\MN{D}$ to $\MN{D}$.
    The left multiplicative identity
    (Theorem~\ref{thm:left_multiplicative_identity})
    gives $E_C(Y^\dagger Y) = E_C(Y^\dagger) E_C(Y)$, so
    $X'^\dagger X'$ is a positive semi-definite fixed point of
    $\E_{B'}$, positive definite by injectivity.
    Hence $X'$ is invertible.

    \emph{Step~5 (Skolem--Noether conclusion).}
    The map $M \mapsto X' M (X')^{-1}$ is a $\C$-algebra automorphism
    of $\MN{D}$ sending $B'^i$ to $A'^i$.
    By Skolem--Noether (Theorem~\ref{thm:skolem_noether}) this is an
    inner automorphism. Undoing the separate gauges gives
    $B^i = \bar{\mu} X A^i X^{-1}$ for all~$i$,
    establishing gauge-phase equivalence
    (Definition~\ref{def:gauge_phase_equiv}).

    The gauging is applied to the individual transfer maps
    $\E_A,\E_B$ separately; the mixed transfer map $F_{AB}$ need not be
    simultaneously unital and trace-preserving.
\end{proof}

\begin{lemma}[Gauged peripheral eigenvector yields Kraus intertwining]%
    \label{thm:gauged_intertwining_core}
    \lean{MPSTensor.gauged_intertwining_core}
    \leanok
    \uses{def:mixed_transfer, thm:qpf, thm:right_canonical_gauge}
    Let $A$ and $B$ be normalized tensors, and let $X \neq 0$ satisfy
    $F_{AB}(X) = \mu X$ with $|\mu| = 1$.
    Choose positive-definite fixed points $\rho_A,\rho_B$ of the transfer maps
    and gauge to unital families
    $A'^i = S_A^{-1} A^i S_A$ and $B'^i = S_B^{-1} B^i S_B$ with
    $S_A S_A^\dagger = \rho_A$ and $S_B S_B^\dagger = \rho_B$.
    Then the transported matrix $X' = S_A^{-1} X (S_B^\dagger)^{-1}$ is nonzero,
    the families $A'$ and $B'$ are unital, and for every~$i$ one has
    \[
        X' (B'^i)^\dagger = \mu (A'^i)^\dagger X',
        \qquad
        A'^i X' = \mu X' B'^i.
    \]
\end{lemma}

\begin{proof}\leanok
    \uses{thm:ks_peripheral, thm:kraus_commute}
    Block-embed the gauged families and transported eigenvector into a
    single unital Kraus map, apply Kadison--Schwarz equality for the
    modulus-one eigenvalue, then extract the intertwining identities
    from the block structure.
\end{proof}

\begin{lemma}[Intertwining yields gauge-phase equivalence]%
    \label{thm:gauged_intertwining_gauge_phase}
    \lean{MPSTensor.gaugePhaseEquiv_of_gauged_intertwining}
    \leanok
    \uses{def:gauge_phase_equiv}
    Let $A$ and $B$ be tensors of the same bond dimension.
    Suppose there are invertible gauges taking them to families $A'$ and $B'$,
    an invertible matrix $X'$, and a scalar $\mu$ with $|\mu| = 1$ such that,
    for every physical index~$i$,
    \[
        A'^i X' = \mu X' B'^i.
    \]
    Then $A$ and $B$ are gauge-phase equivalent.
\end{lemma}

\begin{proof}\leanok
    Compose the two gauge matrices with the intertwiner to obtain a
    single invertible matrix conjugating $A$ to a scalar multiple of~$B$.
\end{proof}

\begin{remark}[Two gauge conventions]\label{rem:gauge_conventions}\leanok
    The right-canonical gauge used here, $A'^i = \rho^{-1/2} A^i \rho^{1/2}$,
    makes $\{A'^i\}$ a \emph{unital} Kraus family ($\sum_i A'^i (A'^i)^\dagger = \Id$).
    A different convention, $\tilde{A}^i = \rho^{1/2} A^i \rho^{-1/2}$, makes
    the family \emph{trace-preserving} ($\sum_i (\tilde{A}^i)^\dagger \tilde{A}^i = \Id$).
    Both conventions appear in practice:
    the unital version via direct matrix square-root inversion
    (used in the transfer-operator gap proof above), and the trace-preserving
    version (used in the canonical-form reduction of
    Chapter~\ref{ch:canonical}).
\end{remark}

\begin{theorem}[Strict transfer-operator gap]\label{thm:transfer_operator_gap}
    \lean{MPSTensor.spectralRadius_mixedTransfer_lt_one}
    \leanok
    \uses{def:gauge_phase_equiv, def:injective, def:mixed_transfer}
    Let $A, B$ be injective normalized MPS tensors that
    are \emph{not} gauge-phase equivalent.
    Then $\rho_{\spec}(F_{AB}) < 1$.
\end{theorem}

\begin{proof}\leanok\uses{thm:spectral_radius_le_one, thm:modulus_one_gauge}
    $\rho_{\spec}(F_{AB}) \le 1$ by
    Lemma~\ref{thm:spectral_radius_le_one}.
    If $\rho_{\spec}(F_{AB}) = 1$,
    Theorem~\ref{thm:modulus_one_gauge} gives gauge-phase
    equivalence, contradicting the hypothesis.
\end{proof}

\section{Transfer-operator gap --- rectangular case}

\begin{lemma}[Rectangular intertwining forces equal bond dimension]%
    \label{thm:rectangular_intertwining_dim_eq}
    \lean{MPSTensor.dim_eq_of_gauged_intertwining}
    \leanok
    \uses{def:injective}
    Let $A$ and $B$ be injective tensors of bond dimensions $D_1$ and $D_2$.
    If there exist $X \in M_{D_1 \times D_2}(\C) \setminus \{0\}$ and
    $\mu \in \C$ such that, for every physical index~$i$,
    \[
        X (B^i)^\dagger = \mu (A^i)^\dagger X,
        \qquad
        A^i X = \mu X B^i,
    \]
    then $D_1 = D_2$.
\end{lemma}

\begin{proof}\leanok
    \uses{def:injective}
    The intertwining relation shows that $\ker(X)$ is invariant under
    all generators of~$B$.
    Since $B$ is injective, its generators span the full matrix algebra,
    so $\ker(X)$ is trivial and $D_2 \le D_1$.
    Applying the same argument to $X^\dagger$ gives $D_1 \le D_2$.
\end{proof}

\begin{theorem}[Rectangular transfer-operator gap]\label{thm:transfer_operator_gap_rect}
    \lean{MPSTensor.mixedTransferSpectralRadius₂_lt_one_of_dim_ne}
    \leanok
    \uses{def:mixed_transfer_rect, def:injective}
    Let $A$ be an injective normalized tensor of bond dimension
    $D_1$ and $B$ an injective normalized tensor of bond
    dimension $D_2$, with $D_1 \neq D_2$.
    Then $\rho_{\spec}(F_{AB}^{\mathrm{rect}}) < 1$.
\end{theorem}

\begin{proof}
    \leanok
    \uses{thm:eigenvalue_bound_rect, thm:gauged_intertwining_core,
        thm:rectangular_intertwining_dim_eq}
    The same word-expansion argument as in
    Theorem~\ref{thm:eigenvalue_bound_rect} gives $|\mu| \le 1$
    for every eigenvalue of $F_{AB}^{\mathrm{rect}}$.
    If $\rho_{\spec}(F_{AB}^{\mathrm{rect}}) = 1$, choose a
    modulus-one eigenvector $X \in M_{D_1 \times D_2}(\C) \setminus \{0\}$.

    Gauge $A,B$ to unital families $A',B'$ and transport $X$ to a nonzero
    rectangular intertwiner $X'$ with $A'^i X' = \mu X' B'^i$ for every~$i$.
    This is the hypothesis of
    Lemma~\ref{thm:rectangular_intertwining_dim_eq}, so $D_1 = D_2$,
    contradicting the hypothesis.
\end{proof}

\begin{theorem}[Rectangular overlap decay]\label{thm:overlap_decay_rect}
    \lean{MPSTensor.mpvOverlap_tendsto_zero_of_dim_ne}
    \leanok
    \uses{def:mpv_overlap, def:injective}
    If $D_1 \neq D_2$ and both tensors are injective and
    normalized, then $O_{AB}(N) \to 0$ as $N \to \infty$.
\end{theorem}

\begin{proof}\leanok\uses{thm:transfer_operator_gap_rect, thm:overlap_trace_rect}
    The transfer-operator gap $\rho_{\spec}(F_{AB}^{\mathrm{rect}}) < 1$
    (Theorem~\ref{thm:transfer_operator_gap_rect}) implies
    $(F_{AB}^{\mathrm{rect}})^N \to 0$, so the operator
    trace converges to~$0$.
    By Lemma~\ref{thm:overlap_trace_rect},
    $O_{AB}(N) = \Tr((F_{AB}^{\mathrm{rect}})^N) \to 0$.
\end{proof}

\section{MPV overlap decay}

\begin{theorem}[Transfer powers converge to zero]\label{thm:transfer_pow_zero}
    \lean{MPSTensor.mixedTransfer_pow_tendsto_zero}
    \leanok
    \uses{def:mixed_transfer, def:gauge_phase_equiv, def:injective}
    Let $A, B$ be injective normalized tensors that are not
    gauge-phase equivalent.
    Then $F_{AB}^n(X) \to 0$ for every $X \in \MN{D}$.
\end{theorem}

\begin{proof}
    \leanok
    \uses{thm:transfer_operator_gap}
    The transfer-operator gap $\rho_{\spec}(F_{AB}) < 1$
    (Theorem~\ref{thm:transfer_operator_gap})
    implies $\|F_{AB}^n\|_{\mathrm{op}} \to 0$
    by the Gelfand formula,
    so $F_{AB}^n(X) \to 0$ for every~$X$.
\end{proof}

\begin{theorem}[Overlap decay]\label{thm:overlap_decay}
    \lean{MPSTensor.mpvOverlap_tendsto_zero}
    \leanok
    \uses{def:mpv_overlap, def:gauge_phase_equiv, def:injective}
    Let $A, B$ be injective normalized tensors of the same bond
    dimension that are not gauge-phase equivalent.
    Then $O_{AB}(N) \to 0$ as $N \to \infty$.
\end{theorem}

\begin{proof}
    \leanok
    \uses{thm:transfer_pow_zero, thm:overlap_trace, lem:trace_expansion}
    By Lemma~\ref{thm:overlap_trace},
    $O_{AB}(N) = \Tr(F_{AB}^N)$.
    Expand the operator trace over matrix units
    (Lemma~\ref{lem:trace_expansion}):
    $\Tr(F_{AB}^N) = \sum_{p,q} (F_{AB}^N(E_{pq}))_{pq}$.
    By Theorem~\ref{thm:transfer_pow_zero},
    each $F_{AB}^N(E_{pq}) \to 0$ entry-wise.
    Since the sum is finite, $O_{AB}(N) \to 0$.
\end{proof}

\section{Block separation}

The block-separation results below record the overlap asymptotics used in the
canonical-form separation of Chapter~\ref{ch:canonical}.

\begin{theorem}[Cross-correlation decay]\label{thm:cross_decay}
    \lean{MPSTensor.cross_correlation_tendsto_zero}
    \leanok
    \uses{def:mixed_transfer, def:gauge_phase_equiv, def:injective}
    If $A$ and $B$ are injective normalized tensors that are not
    gauge-phase equivalent, then for every $X \in \MN{D}$ one has
    $\tr(F_{AB}^N(X)) \to 0$ as $N \to \infty$.
\end{theorem}

\begin{proof}\leanok\uses{thm:transfer_pow_zero}
    By Theorem~\ref{thm:transfer_pow_zero}, $F_{AB}^N(X) \to 0$ in operator
    norm, so $\tr(F_{AB}^N(X)) \to 0$.
\end{proof}

\begin{theorem}[Self-correlation persists]\label{thm:self_persist}
    \lean{MPSTensor.self_correlation_persists}
    \leanok
    \uses{def:transfer_map}
    If $\rho$ is a fixed point of $\E_A$, then
    $\tr(\E_A^N(\rho)) = \tr(\rho)$ for all $N \ge 0$.
    In particular this applies to the distinguished QPF fixed point from
    Theorem~\ref{thm:qpf}.
\end{theorem}

\begin{proof}\leanok
    Since $\rho$ is a fixed point, $\E_A^N(\rho) = \rho$ for every $N$, so
    $\tr(\E_A^N(\rho)) = \tr(\rho)$.
\end{proof}

\section{Transfer-operator gap under irreducible-TP hypotheses}

The preceding transfer-operator gap results require \emph{injectivity} of both
tensors. The following variants weaken this to \emph{irreducibility}
combined with the trace-preserving normalization~(\ref{eq:spec_tp_normalization}), following Cirac et
al.~\cite{Cirac2017Annals} and the irreducibility theory of
\cite[Section~6.2]{Wolf2012Quantum}. The argument replaces the
Kraus-commutation step with a Cauchy--Schwarz rigidity argument for
the Perron--Frobenius fixed point.

\begin{theorem}[Gauge-equivalence from irreducible-TP spectral
        radius]\label{thm:modulus_one_gauge_irr_tp}
    \lean{MPSTensor.modulus_one_eigenvalue_implies_gauge_of_irreducible_TP}
    \leanok
    \uses{def:mixed_transfer, def:gauge_phase_equiv,
        def:irreducible_tensor, thm:eigenvalue_bound}
    Let $A, B$ be irreducible trace-preserving normalized MPS tensors
    of bond dimension~$D$.
    If $\rho_{\spec}(F_{AB}) \ge 1$, then $A$ and $B$ are gauge-phase equivalent.
\end{theorem}

\begin{proof}
    \leanok
    \uses{thm:modulus_one_gauge, thm:qpf, thm:right_canonical_gauge,
        thm:ks_peripheral, thm:kraus_commute,
        thm:left_multiplicative_identity, thm:psd_fp_pd_irred,
        thm:gauged_intertwining_core, thm:gauged_intertwining_gauge_phase,
        thm:skolem_noether}
    Choose a modulus-one eigenvector as in
    Theorem~\ref{thm:modulus_one_gauge}.
    The proof follows the five-step structure of
    Theorem~\ref{thm:modulus_one_gauge}, with one key adaptation.
    Steps~1--3 (separate gauging, block Kraus family, KS equality and
    Kraus commutation) proceed identically: the gauging uses the
    unique positive definite fixed points guaranteed by irreducibility
    (Theorem~\ref{thm:qpf}) and the right-canonical gauge
    (Theorem~\ref{thm:right_canonical_gauge}).
    The change is in Step~4: the Kadison--Schwarz equality gives
    $X'^\dagger X'$ as a nonzero positive semidefinite fixed point of
    $\E_{B'}$.
    Under injectivity this is immediately positive definite (the fixed
    point is unique up to scaling).
    Under the weaker hypothesis of irreducibility,
    Theorem~\ref{thm:psd_fp_pd_irred} upgrades it to positive
    definite, giving invertibility of~$X'$.
    Formally, Lemma~\ref{thm:gauged_intertwining_core} combines the
    separately gauged eigenvector into the required intertwining
    relations.
    Once $X'$ is invertible, Lemma~\ref{thm:gauged_intertwining_gauge_phase}
    upgrades the intertwining to gauge-phase equivalence.
\end{proof}

\begin{theorem}[Strict transfer-operator gap (irreducible-TP)]%
    \label{thm:transfer_operator_gap_irr_tp}
    \lean{MPSTensor.spectralRadius_mixedTransfer_lt_one_of_irreducible_TP}
    \leanok
    \uses{def:gauge_phase_equiv, def:irreducible_tensor,
        def:mixed_transfer}
    Let $A, B$ be irreducible trace-preserving normalized MPS tensors
    of the same bond dimension that are not gauge-phase equivalent.
    Then $\rho_{\spec}(F_{AB}) < 1$.
\end{theorem}

\begin{proof}\leanok\uses{thm:spectral_radius_le_one,
        thm:modulus_one_gauge_irr_tp}
    Contrapositive of
    Theorem~\ref{thm:modulus_one_gauge_irr_tp}, combined with
    $\rho_{\spec}(F_{AB}) \le 1$
    (Lemma~\ref{thm:spectral_radius_le_one}).
\end{proof}

\begin{theorem}[Transfer power decay (irreducible-TP)]%
    \label{thm:transfer_pow_zero_irr_tp}
    \lean{MPSTensor.mixedTransfer_pow_tendsto_zero_of_irreducible_TP}
    \leanok
    \uses{def:mixed_transfer, def:gauge_phase_equiv,
        def:irreducible_tensor}
    Under the hypotheses of
    Theorem~\ref{thm:transfer_operator_gap_irr_tp},
    $F_{AB}^n(X) \to 0$ for every $X \in \MN{D}$.
\end{theorem}

\begin{proof}\leanok\uses{thm:transfer_operator_gap_irr_tp}
    By the Gelfand formula,
    $\rho_{\spec}(F_{AB}) < 1$ implies $\|F_{AB}^n\|_{\mathrm{op}} \to 0$.
\end{proof}

\begin{theorem}[Overlap decay (irreducible-TP)]%
    \label{thm:overlap_decay_irr_tp}
    \lean{MPSTensor.mpvOverlap_tendsto_zero_of_irreducible_TP}
    \leanok
    \uses{def:mpv_overlap, def:gauge_phase_equiv,
        def:irreducible_tensor}
    Let $A, B$ be irreducible trace-preserving normalized tensors of
    the same bond dimension that are not gauge-phase equivalent.
    Then $O_{AB}(N) \to 0$ as $N \to \infty$.
\end{theorem}

\begin{proof}\leanok\uses{thm:transfer_operator_gap_irr_tp, thm:overlap_trace}
    Combine the transfer-operator gap
    (Theorem~\ref{thm:transfer_operator_gap_irr_tp})
    with the overlap--trace identity
    (Lemma~\ref{thm:overlap_trace}).
\end{proof}

\begin{theorem}[Rectangular transfer-operator gap (irreducible-TP)]%
    \label{thm:transfer_operator_gap_rect_irr_tp}
    \lean{MPSTensor.mixedTransferSpectralRadius₂_lt_one_of_dim_ne_of_irreducible_TP}
    \leanok
    \uses{def:mixed_transfer_rect, def:irreducible_tensor}
    Let $A$ (bond dimension $D_1$) and $B$ (bond dimension $D_2$) be
    irreducible trace-preserving normalized tensors with
    $D_1 \neq D_2$.
    Then $\rho_{\spec}(F_{AB}^{\mathrm{rect}}) < 1$.
\end{theorem}

\begin{proof}
    \leanok
    \uses{thm:eigenvalue_bound, thm:psd_fp_pd_irred}
    Assume for contradiction that $\rho_{\spec}(F_{AB}^{\mathrm{rect}}) = 1$.
    Choose a modulus-one eigenvector $X \in M_{D_1 \times D_2}(\C)$.
    Gauge $A,B$ separately to unital families $A',B'$ and transport $X$ to
    $X'\neq 0$.
    The matrix $X'X'^\dagger$ ($D_1 \times D_1$) is a nonzero
    positive-semidefinite fixed point of $\E_{A'}$, so by
    Theorem~\ref{thm:psd_fp_pd_irred} it is positive definite, hence invertible;
    thus $X'$ has full row rank and $D_1 \le D_2$. Symmetrically
    $X'^\dagger X'$ ($D_2 \times D_2$) is positive definite, giving full column
    rank and $D_2 \le D_1$. Hence $D_1=D_2$, contradicting the hypothesis.
\end{proof}

\begin{theorem}[Rectangular overlap decay (irreducible-TP)]%
    \label{thm:overlap_decay_rect_irr_tp}
    \lean{MPSTensor.mpvOverlap_tendsto_zero_of_dim_ne_of_irreducible_TP}
    \leanok
    \uses{def:mpv_overlap, def:irreducible_tensor}
    If $D_1 \neq D_2$ and both tensors are irreducible and
    trace-preserving normalized, then
    $O_{AB}(N) \to 0$ as $N \to \infty$.
\end{theorem}

\begin{proof}\leanok\uses{thm:transfer_operator_gap_rect_irr_tp,
        thm:overlap_trace_rect}
    Combine
    Theorem~\ref{thm:transfer_operator_gap_rect_irr_tp} with
    Lemma~\ref{thm:overlap_trace_rect}.
\end{proof}

\section{Conditional expectation from a faithful fixed point}

This section follows~\cite[Section~6.4]{Wolf2012Quantum}; the conditional
expectation is built from the fixed-point algebra of~\cite[Theorem~6.14]{Wolf2012Quantum}
via~\cite[(1.40)]{Wolf2012Quantum}.
The scalar conditional expectation $E_\sigma$ is the special case where
the fixed-point $*$-subalgebra is the scalar algebra $\C\cdot\Id$
(the primitive-channel case). The general irreducible case with
nontrivial period requires the Wedderburn block decomposition
of~\cite[Theorem~6.14]{Wolf2012Quantum}.
% The nontrivial-period conditional expectation is outside the present
% formal scope of this section.

\begin{definition}[Conditional expectation]\label{def:is_conditional_expectation}
    \lean{Kraus.IsConditionalExpectation}
    \leanok
    Given a $*$-subalgebra $S \subseteq \MN{D}$, a linear map
    $P : \MN{D} \to \MN{D}$ is a conditional expectation onto $S$ if it is
    idempotent, unital, has range contained in $S$, and satisfies
    $P(X)=X$ for all $X \in S$.
\end{definition}

\begin{definition}[Scalar conditional expectation]\label{def:scalar_conditional_expectation}
    \lean{Kraus.scalarConditionalExpectation}
    \leanok
    \uses{def:is_conditional_expectation}
    For $\sigma \ge 0$ with $\tr(\sigma)\neq 0$, define
    \[
        E_\sigma(X) := \frac{\tr(\sigma X)}{\tr(\sigma)} \Id.
    \]
\end{definition}

\begin{lemma}[Evaluation formula]\label{lem:scalar_conditional_expectation_apply}
    \lean{Kraus.scalarConditionalExpectation_apply}
    \leanok
    \uses{def:scalar_conditional_expectation}
    $E_\sigma(X)=\dfrac{\tr(\sigma X)}{\tr(\sigma)} \Id$.
\end{lemma}

\begin{lemma}[Unitality]\label{thm:scalar_conditional_expectation_unital}
    \lean{Kraus.scalarConditionalExpectation_unital}
    \leanok
    \uses{def:scalar_conditional_expectation}
    $E_\sigma(\Id)=\Id$.
\end{lemma}

\begin{lemma}[Idempotence]\label{thm:scalar_conditional_expectation_idempotent}
    \lean{Kraus.scalarConditionalExpectation_idempotent}
    \leanok
    \uses{def:scalar_conditional_expectation}
    $E_\sigma(E_\sigma(X)) = E_\sigma(X)$ for all $X$.
\end{lemma}

\begin{lemma}[Range is scalar]\label{lem:scalar_conditional_expectation_range_scalar}
    \lean{Kraus.scalarConditionalExpectation_range_scalar}
    \leanok
    \uses{def:scalar_conditional_expectation}
    For every $X$, there exists $c\in\C$ with $E_\sigma(X)=c\,\Id$.
\end{lemma}

\begin{lemma}[Fixes scalar operators]\label{lem:scalar_conditional_expectation_fixes_scalar}
    \lean{Kraus.scalarConditionalExpectation_fixes_scalar}
    \leanok
    \uses{def:scalar_conditional_expectation}
    For $c\in\C$, $E_\sigma(c\,\Id)=c\,\Id$.
\end{lemma}

\begin{lemma}[Absorbs the adjoint map]\label{thm:scalar_conditional_expectation_absorbs_adjointMap}
    \lean{Kraus.scalarConditionalExpectation_absorbs_adjointMap}
    \leanok
    \uses{def:scalar_conditional_expectation}
    If $T(\sigma)=\sigma$, then
    $E_\sigma\circ T^* = E_\sigma$.
\end{lemma}

\begin{lemma}[Adjoint map absorbs scalar projection]
    \label{thm:adjointMap_absorbs_scalarConditionalExpectation}
    \lean{Kraus.adjointMap_absorbs_scalarConditionalExpectation}
    \leanok
    \uses{def:scalar_conditional_expectation, def:tp_kraus}
    If $T$ is trace-preserving, then
    $T^*\circ E_\sigma = E_\sigma$.
\end{lemma}

\begin{lemma}[Image lies in adjoint fixed points]
    \label{lem:scalar_conditional_expectation_mem_adjointFixedPoints}
    \lean{Kraus.scalarConditionalExpectation_mem_adjointFixedPoints}
    \leanok
    \uses{thm:adjointMap_absorbs_scalarConditionalExpectation}
    For all $X$, $E_\sigma(X)$ is fixed by $T^*$.
\end{lemma}

\begin{theorem}[Scalar map is a conditional expectation]
    \label{thm:scalar_conditional_expectation_isConditionalExpectation}
    \lean{Kraus.scalarConditionalExpectation_isConditionalExpectation}
    \leanok
    \uses{def:is_conditional_expectation, def:scalar_conditional_expectation,
        def:tp_kraus,
        thm:scalar_conditional_expectation_unital,
        thm:scalar_conditional_expectation_idempotent,
        lem:scalar_conditional_expectation_mem_adjointFixedPoints,
        lem:scalar_conditional_expectation_fixes_scalar}
    Let $K$ be trace-preserving, let $\rho>0$ be a positive-definite fixed
    point of the associated channel $T(\rho)=\sum_i K_i\rho K_i^\dagger$ (so $T(\rho)=\rho$), and
    assume every adjoint fixed point of $T^*$ is a scalar multiple of $\Id$.
    Then $E_\rho$ is a conditional expectation onto the adjoint fixed-point
    $*$-subalgebra.
\end{theorem}

\begin{lemma}[Complete positivity of the scalar conditional expectation]
    \label{thm:scalar_conditional_expectation_isCPMap}
    \lean{Kraus.scalarConditionalExpectation_isCPMap}
    \leanok
    \uses{def:scalar_conditional_expectation}
    For any positive semidefinite $\sigma$, the scalar conditional
    expectation $E_\sigma$ is a completely positive map.
\end{lemma}
\begin{proof}\leanok
    Write $S=\sqrt{\sigma}$ and define the Kraus family
    $K_{j,i}=\ketbra{j}{i}\,S$ indexed by $(j,i)\in\{1,\ldots,D\}^2$.
    Using $\ketbra{j}{i}\,M\,\ketbra{i}{j} = M_{ii} \ketbra{j}{j}$ and
    summing first over $i$ and then over $j$,
    \[
        \sum_{j,i} K_{j,i}\,X\,K_{j,i}^\dagger
        = \sum_j \ketbra{j}{j} \tr(SXS)
        = \tr(\sigma X) \Id,
    \]
    where we used cyclicity of the trace and $S^2=\sigma$. Hence
    $E_\sigma = (\tr(\sigma))^{-1}\cdot(X\mapsto\sum_{j,i}K_{j,i}XK_{j,i}^\dagger)$
    is the Kraus map scaled by the non-negative real number
    $(\tr(\sigma))^{-1}$, and therefore completely positive.
\end{proof}

\begin{lemma}[$\sigma$-trace preservation of the scalar conditional expectation]
    \label{thm:scalar_conditional_expectation_preserves_weighted_trace}
    \lean{Kraus.scalarConditionalExpectation_preserves_weighted_trace}
    \leanok
    \uses{def:scalar_conditional_expectation}
    If $\tr(\sigma)\neq 0$, then for every $X$
    \[
        \tr(\sigma\,E_\sigma(X))
        = \tr(\sigma X).
    \]
\end{lemma}
\begin{proof}\leanok
    Direct computation:
    $\tr(\sigma\cdot\tfrac{\tr(\sigma X)}{\tr(\sigma)}\Id)
        = \tfrac{\tr(\sigma X)}{\tr(\sigma)}\tr(\sigma)
        = \tr(\sigma X)$.
\end{proof}

\section{Stationary support}

This section develops the support-projection theory
behind stationary states, following~\cite[Section~6.4, Propositions~6.10--6.11]{Wolf2012Quantum}.

\begin{lemma}[Lower-triangular Kraus condition implies corner invariance]%
    \label{lem:lower_zero_invariance}
    \lean{MPSTensor.lowerZero_implies_invariance}
    \leanok
    Let $A$ be an MPS tensor and let $P$ be an orthogonal projection.
    If, for every physical index~$i$,
    \[
        (\Id - P) A^i P = 0
    \]
    then
    \[
        P \E_A(P X P) P = \E_A(P X P)
    \]
    for all $X \in \MN{D}$.
\end{lemma}

\begin{proof}\leanok
    Expand the transfer map and use $(\Id - P) A^i P = 0$ to reduce
    each Kraus term to the $P$-corner.
\end{proof}

\begin{lemma}[Support projection invariance]%
    \label{thm:support_proj_fixed}
    \lean{Channel.support_proj_fixed}
    \leanok
    \uses{def:channel}
    Let $E$ be a quantum channel and $\rho \ge 0$ a positive
    semidefinite fixed point ($E(\rho) = \rho$).
    Let $P$ denote the support projection of~$\rho$.
    Then
    \[
        P E(P X P) P = E(P X P)
    \]
    for all $X \in \MN{D}$.
\end{lemma}

\begin{proof}
    \leanok
    \uses{lem:lower_zero_invariance}
    Write $E$ in Kraus form $E(X) = \sum_i K_i X K_i^\dagger$.
    From $E(\rho) = \rho$ and $\rho \ge 0$ one deduces
    $(\Id - P) K_i P = 0$ for all~$i$.
    Apply Lemma~\ref{lem:lower_zero_invariance}.
\end{proof}

\begin{definition}[Stationary state]\label{def:stationaryState}
    \lean{Channel.stationaryState}
    \leanok
    \uses{def:channel, def:irreducible_cp}
    For an irreducible channel $E$ on $\MN{D}$ (with $D \ge 1$),
    the \emph{stationary state} is the unique density-matrix fixed
    point $\rho_\infty$ satisfying $E(\rho_\infty) = \rho_\infty$,
    $\rho_\infty > 0$, and $\tr(\rho_\infty) = 1$.
    Uniqueness follows from the quantum Perron--Frobenius theorem
    (Theorem~\ref{thm:qpf}).
\end{definition}

\begin{definition}[Stationary support]\label{def:stationarySupport}
    \lean{Channel.stationarySupport}
    \leanok
    \uses{def:stationaryState}
    The \emph{stationary support} of an irreducible channel~$E$
    is the support projection of the stationary state
    (Definition~\ref{def:stationaryState}).
\end{definition}

\begin{theorem}[Stationary support is full]%
    \label{thm:stationarySupport_eq_one}
    \lean{Channel.stationarySupport_eq_one}
    \leanok
    \uses{def:stationarySupport, def:irreducible_cp}
    For an irreducible channel~$E$, the stationary support equals
    the identity: $\mathrm{supp}(\rho_\infty) = \Id$.
\end{theorem}

\begin{proof}
    \leanok
    \uses{thm:support_proj_fixed}
    The support projection $P$ of $\rho_\infty$ is an orthogonal
    projection satisfying $P E(PXP) P = E(PXP)$ for all~$X$
    (Lemma~\ref{thm:support_proj_fixed}).
    By irreducibility, $P = 0$ or $P = \Id$.
    Since $\rho_\infty \neq 0$, we have $P \neq 0$, so $P = \Id$.
\end{proof}

\begin{remark}[Stationary-support formulations]%
    \label{rem:stationary_support_todo}
    The non-trivial formulations needed here are the following results from
    Wolf:
    \begin{itemize}
        \item the converse direction of~\cite[Proposition~6.10]{Wolf2012Quantum}:
              if the support projection of every positive semidefinite fixed point
              is full, then the channel is irreducible;
        \item \cite[Proposition~6.11]{Wolf2012Quantum}: minimality of the
              stationary support among invariant projections, without assuming
              irreducibility.
    \end{itemize}
    These statements are not used elsewhere in this chapter.
\end{remark}

\subsection{Faithful compression onto the support sector}

This subsection establishes the results on the corner algebra,
$*$-structure, and compression of a PSD fixed point onto its support sector
needed for~\cite[Corollary~6.6]{Wolf2012Quantum}: the corner carries a
$\C$-algebra and $*$-structure, the two standard presentations of the corner
are identified, and the compression of a PSD fixed point onto its support
sector is positive definite.

\begin{lemma}[Support projection annihilates the kernel]%
    \label{lem:supportProj_mulVec_eq_zero_of_mulVec_eq_zero}
    \lean{MPSTensor.supportProj_mulVec_eq_zero_of_mulVec_eq_zero}
    \leanok
    Let $\rho \succeq 0$ on $\C^D$ with support projection
    $P = \mathrm{supp}(\rho)$. For every $v \in \C^D$, if $\rho\, v = 0$,
    then $P\, v = 0$.
\end{lemma}
\begin{proof}\leanok
    Diagonalize $\rho = U\,\diag(\lambda)\,U^\dagger$ with
    $\lambda_j \ge 0$. Writing $w := U^\dagger v$, the hypothesis
    $\rho v = 0$ gives $\lambda_j\, w_j = 0$ for every $j$; hence $w_j = 0$
    whenever $\lambda_j > 0$. The support projection acts as
    $\diag(\chi_{\lambda_j > 0})$ in the eigenbasis, so
    $P\,v = U\,\diag(\chi_{\lambda_j > 0})\,w = 0$.
\end{proof}

\begin{definition}[Corner submodule / idempotent-corner identification]%
    \label{def:cornerSubmoduleCornerEquiv}
    \lean{MatrixCorner.cornerSubmoduleCornerEquiv}
    \leanok
    For an idempotent $P \in \MN{D}$, the subspace
    $\{X \mid P X P = X\}$ and the corner
    $\{P\,X\,P \mid X \in \MN{D}\}$ are identified as $\C$-linear spaces by
    the identity on underlying matrices.
\end{definition}

\begin{remark}[$\C$-algebra and $*$-structure on the corner]%
    \label{rem:corner_algebra_star_instances}
    \lean{MatrixCorner.instAlgebraComplexCorner}
    \lean{MatrixCorner.cornerStar}
    \lean{MatrixCorner.cornerStarRing}
    \lean{MatrixCorner.cornerStarModuleComplex}
    \leanok
    The corner algebra associated to an idempotent already carries a ring
    structure with unit~$P$. In the matrix case it also has
    $\C$-module and $\C$-algebra structures and, under the additional
    assumption $P^\dagger = P$, star, star ring, and star module
    structures over~$\C$.
\end{remark}

\begin{theorem}[Faithful compression onto the support]%
    \label{thm:compression_on_support_posDef}
    \lean{Matrix.PosSemidef.compression_on_support_posDef}
    \leanok
    \uses{lem:supportProj_mulVec_eq_zero_of_mulVec_eq_zero}
    Let $\rho \succeq 0$ on $\C^D$ with support projection
    $P = \mathrm{supp}(\rho)$, and let
    $V : \C^D \to \C^k$ be a compression isometry satisfying
    $V\,V^\dagger = \Id_k$ and $V^\dagger\,V = P$. Then the compression
    \[
        V\,\rho\,V^\dagger \in \MN{k}
    \]
    is positive definite.
\end{theorem}
\begin{proof}\leanok
    \uses{lem:supportProj_mulVec_eq_zero_of_mulVec_eq_zero}
    Hermiticity follows from Hermiticity of $\rho$ and the product-star
    identity. For strict positivity, let $w \neq 0$, set
    $u := V^\dagger w$, and note that $V V^\dagger = \Id_k$ forces
    $u \neq 0$, while $V^\dagger V = P$ forces $P u = u$. The quadratic
    form computes as
    $w^\dagger (V\rho V^\dagger) w = u^\dagger \rho u$,
    which is non-negative by positive semidefiniteness. If it vanishes, then
    $\rho u = 0$, and by
    Lemma~\ref{lem:supportProj_mulVec_eq_zero_of_mulVec_eq_zero} also
    $P u = 0$, contradicting $P u = u \neq 0$.
\end{proof}

\section{Wedderburn decomposition of the fixed-point algebra}

This section proves the structure theorem for the fixed-point
algebra of a trace-preserving Kraus map, following~\cite[Theorems~6.12--6.14]{Wolf2012Quantum}. The proof uses the
Jacobson radical characterization of semisimplicity together with the
Wedderburn--Artin structure theorem for finite-dimensional semisimple
algebras over an algebraically closed field.

\begin{theorem}[Fixed-point algebra is semisimple]%
    \label{thm:fixedPointAlgebra_isSemisimpleRing}
    \lean{Kraus.fixedPointAlgebra_isSemisimpleRing}
    \leanok
    \uses{thm:adjointFixedPointsStarSubalgebra}
    The adjoint-fixed-point $*$-subalgebra of a trace-preserving
    Kraus map on $\MN{D}$ is a semisimple ring.
\end{theorem}

\begin{proof}
    \leanok
    \uses{thm:adjointFixedPointsStarSubalgebra}
    The subalgebra is finite-dimensional (as a subalgebra of~$\MN{D}$),
    hence Artinian. Following~\cite[Theorem~6.14]{Wolf2012Quantum}, it
    suffices to show the Jacobson radical~$J$ vanishes. If $x \in J$, then
    $x^*x \in J$
    (left-ideal closure). The radical of an Artinian ring is nilpotent,
    so $x^*x$ is nilpotent. A self-adjoint nilpotent matrix satisfies
    $\lVert x^*x\rVert^{2^k} = \lVert (x^*x)^{2^k}\rVert = 0$,
    hence $x^*x = 0$, and the C$^*$-identity gives $x = 0$.
\end{proof}

\begin{theorem}[Abstract Wedderburn--Artin decomposition]%
    \label{thm:fixedPointAlgebra_wedderburnArtin}
    \lean{Kraus.fixedPointAlgebra_wedderburnArtin}
    \leanok
    \uses{thm:fixedPointAlgebra_isSemisimpleRing}
    There exist $n\in\N$ and dimensions $d_1,\ldots,d_n\ge 1$ such that
    the adjoint-fixed-point $*$-subalgebra is $\C$-algebra isomorphic to
    \[
        \prod_{k=1}^{n} \MN{d_k}.
    \]
\end{theorem}

\begin{proof}
    \leanok
    \uses{thm:fixedPointAlgebra_isSemisimpleRing}
    Combine semisimplicity (Theorem~\ref{thm:fixedPointAlgebra_isSemisimpleRing})
    with the Wedderburn--Artin structure theorem for finite-dimensional
    semisimple algebras over an algebraically closed field.
\end{proof}

\begin{theorem}[Fixed-point algebra decomposition with multiplicities]%
    \label{thm:adjointFixedPoints_wedderburnDecomp}
    \lean{Kraus.adjointFixedPoints_wedderburnDecomp}
    \leanok
    \uses{thm:fixedPointAlgebra_wedderburnArtin}
    There exist integers $n$, block sizes $d_1,\ldots,d_n \ge 1$,
    multiplicities $m_1,\ldots,m_n \ge 1$, and a $\C$-algebra isomorphism from
    the adjoint-fixed-point $*$-subalgebra to
    \[
        \prod_{k=1}^{n} \MN{d_k}
    \]
    such that
    \[
        \sum_{k=1}^{n} d_k m_k \le D.
    \]
\end{theorem}
\begin{proof}
    \leanok
    \uses{thm:fixedPointAlgebra_wedderburnArtin}
    Apply the abstract Wedderburn decomposition
    (Theorem~\ref{thm:fixedPointAlgebra_wedderburnArtin}).
    Since the adjoint-fixed-point algebra is realized as a subalgebra of the
    ambient matrix algebra $\MN{D}$, this realization also yields
    multiplicities $m_k$ and the bound $\sum_k d_k m_k \le D$.
\end{proof}

\begin{remark}[Block form of the fixed-point algebra]
    Wolf's block form identifies the adjoint-fixed-point algebra, up to a
    unitary change of basis, with
    \[
        0 \oplus \bigoplus_k \MN{d_k} \otimes \Id_{m_k}
    \]
    together with the density-block form $\MN{d_k} \otimes \rho_k$.
    The theorem above gives the product decomposition, multiplicities, and
    dimension bound needed from this block form.
\end{remark}

\begin{theorem}[Weighted fixed points form a $*$-subalgebra]%
    \label{thm:weightedFixedPointsStarSubalgebra}
    \lean{Kraus.weightedFixedPointsStarSubalgebra}
    \leanok
    \uses{def:tp_kraus, thm:fixedPointsStarSubalgebra}
    Let $T(X) = \sum_i K_i X K_i^\dagger$ be trace-preserving, and let
    $\rho > 0$ satisfy $T(\rho) = \rho$. If
    \[
        F_T := \{X \in \MN{D} \mid T(X) = X\},
    \]
    then $\rho^{-1/2} F_T \rho^{-1/2}$ is a $*$-subalgebra of $\MN{D}$.
\end{theorem}
\begin{proof}
    \leanok
    \uses{thm:fixedPointsStarSubalgebra}
    Set $B_i = \rho^{-1/2} K_i \rho^{1/2}$. The equation $T(\rho) = \rho$
    implies that the family $\{B_i\}$ is unital, and trace preservation of $T$
    implies that $\rho$ is a positive-definite fixed point of the adjoint map
    of the gauged family. Apply Theorem~\ref{thm:fixedPointsStarSubalgebra}
    to $B$. The identity
    \[
        \sum_i B_i X B_i^\dagger = X
        \iff
        T(\rho^{1/2} X \rho^{1/2}) = \rho^{1/2} X \rho^{1/2}
    \]
    identifies the fixed points of the gauged map with
    $\rho^{-1/2} F_T \rho^{-1/2}$.
\end{proof}

\section{Further irreducibility and primitivity equivalences}

This section collects several criteria and equivalences from~\cite[Theorems~6.2, 6.7, 6.8]{Wolf2012Quantum} that are
used later.

\begin{theorem}[Exponential positivity condition]%
    \label{thm:wolf_6_2_item_3}
    \lean{exp_posDef_of_irreducible_cp}
    \lean{irreducible_iff_exp_posDef_forall}
    \leanok
    \uses{def:irreducible_cp}
    Let $E$ be an irreducible completely positive map, $A \ge 0$ with
    $A \neq 0$, and $t > 0$. Then $\exp(tE)(A)$ is positive definite.
    Conversely, if $\exp(tE)(A)$ is positive definite for every such $t$, $A$,
    then $E$ is irreducible (full equivalence \lean{irreducible_iff_exp_posDef_forall}).
\end{theorem}

\begin{theorem}[Kraus-span equivalence]%
    \label{thm:wolf_6_8_kraus_span}
    \lean{MPSTensor.wolf_theorem_6_8_kraus_span}
    \leanok
    \uses{def:primitive_channel}
    Under normalization, primitivity in Wolf's sense is equivalent to eventual
    full Kraus rank.
\end{theorem}

\begin{theorem}[Combined primitivity equivalence]%
    \label{thm:wolf_6_8_conjunction}
    \lean{MPSTensor.wolf_theorem_6_8_conjunction}
    \leanok
    \uses{thm:wolf_6_8_kraus_span}
    Under normalization, primitivity in Wolf's sense is equivalent to the conjunction
    of eventual full Kraus rank, normality, and strong irreducibility.
\end{theorem}

\begin{remark}[Pairwise primitivity equivalences]
    \label{rem:wolf_6_8_pairwise}
    \lean{MPSTensor.primitivePaper_iff_hasEventuallyFullKrausRank}
    \lean{MPSTensor.primitivePaper_iff_stronglyIrreducible}
    \lean{MPSTensor.hasEventuallyFullKrausRank_iff_stronglyIrreducible}
    \lean{MPSTensor.hasEventuallyFullKrausRank_iff_isNormal}
    \leanok
    These are the pairwise equivalences in
    \cite[Theorem~6.8]{Wolf2012Quantum}.
\end{remark}

\begin{theorem}[Scalar fixed-point algebra case]
    \label{thm:wolf_6_15_scalar_conditional_expectation}
    \lean{Kraus.scalarConditionalExpectation_isConditionalExpectation}
    \leanok
    \uses{def:channel}
    In the scalar fixed-point algebra setting, the weighted map
    $X \mapsto \frac{\tr(\rho X)}{\tr(\rho)}\Id$ is a conditional
    expectation onto the adjoint fixed-point $*$-subalgebra.
\end{theorem}

\section{Multi-cycle block-permutation structure}

This section develops the block-permutation structure underlying the
cycle decomposition of~\cite[Theorem~6.16]{Wolf2012Quantum}. The
multi-cycle decomposition refines the single-cycle case to handle
arbitrary peripheral periods.

\begin{definition}[Block-permutation structure]
    \label{def:cycle_structure}
    \lean{CycleStructure}
    \leanok
    Let $T : \MN{D} \to \MN{D}$. A \emph{block-permutation structure} for
    $T$ consists of a finite index set $\iota$, a family of orthogonal
    projections $P_k \in \MN{D}$ for $k \in \iota$, and a permutation
    $\sigma \in \mathrm{Sym}(\iota)$ (not necessarily a single cycle---in
    general $\sigma$ may have multiple disjoint cycles), such that
    \[
        T(P_{\sigma(k)}) = P_k,\quad
        T(P_k X) = T(P_k)\,T(X),\quad
        T(X P_k) = T(X)\,T(P_k)
    \]
    for every $k \in \iota$ and $X \in \MN{D}$.
\end{definition}

\begin{definition}[Multi-cycle block-permutation structure]
    \label{def:multi_cycle_decomposition}
    \lean{MultiCycleDecomposition}
    \leanok
    A \emph{multi-cycle decomposition} of $T : \MN{D} \to \MN{D}$ consists of a
    finite cycle index $\iota$, a per-cycle period
    $m : \iota \to \N_{>0}$, a projection family
    $P_{c,k} \in \MN{D}$ for $c \in \iota$ and $k \in \Fin{m(c)}$
    consisting of orthogonal projections. For every $c \in \iota$ and
    $k \in \Fin{m(c)}$, the per-cycle cyclic action is
    \[
        T(P_{c,k+1}) = P_{c,k}.
    \]
    The multiplicative-domain factorisations are
    $T(P_{c,k}\,X) = T(P_{c,k})\,T(X)$ and
    $T(X\,P_{c,k}) = T(X)\,T(P_{c,k})$.
\end{definition}

\begin{lemma}[Per-cycle corner preservation]
    \label{thm:multi_cycle_preserves_corner_pow_period}
    \lean{MultiCycleDecomposition.preserves_corner_pow_period}
    \leanok
    \uses{def:multi_cycle_decomposition}
    For every cycle $c \in \iota$ and index $k \in \Fin{m(c)}$, the
    iterate $T^{m(c)}$ preserves the corner $P_{c,k} \MN{D}\,P_{c,k}$.
\end{lemma}

\begin{proof}
    \leanok
    \uses{def:multi_cycle_decomposition}
    By induction on $n$, the $n$-th iterate $T^n$ sends
    $P_{c,k+n}\,X\,P_{c,k+n}$ to $P_{c,k}\,T^n(X)\,P_{c,k}$, using the
    multiplicative-domain factorisations and the cyclic action
    $T(P_{c,j+1}) = P_{c,j}$ at each induction step. Specialising to
    $n = m(c)$ and using that the indices $k + m(c) = k$ in $\Fin{m(c)}$
    yields corner preservation of $T^{m(c)}$.
\end{proof}

\begin{lemma}[Common-period corner preservation]
    \label{thm:multi_cycle_preserves_corner_pow_of_dvd}
    \lean{MultiCycleDecomposition.preserves_corner_pow_of_dvd}
    \leanok
    \uses{thm:multi_cycle_preserves_corner_pow_period}
    If $N \in \N$ is divisible by every per-cycle period $m(c)$,
    then $T^N$ preserves every corner $P_{c,k} \MN{D}\,P_{c,k}$ of the
    decomposition.
\end{lemma}

\begin{proof}
    \leanok
    \uses{thm:multi_cycle_preserves_corner_pow_period}
    Write $N = m(c)\,q$. Corner preservation is stable under taking
    powers, so $T^N = (T^{m(c)})^q$ preserves each corner by
    Lemma~\ref{thm:multi_cycle_preserves_corner_pow_period}.
\end{proof}

\begin{definition}[Flattening to a single-index block-permutation structure]
    \label{def:multi_cycle_to_cycle_structure}
    \lean{MultiCycleDecomposition.toCycleStructure}
    \leanok
    \uses{def:multi_cycle_decomposition, def:cycle_structure}
    A multi-cycle decomposition gives a single-index block-permutation
    structure on the sigma index $\Sigma_{c \in \iota} \Fin{m(c)}$: the
    permutation is the sigma-product of per-cycle cyclic shifts
    $k \mapsto k+1$ (whose cycle decomposition has one cycle per
    $c \in \iota$), the projections are $(c, k) \mapsto P_{c,k}$, and the
    multiplicative-domain factorisations descend componentwise. The
    resulting map forgets the explicit cycle-indexing.
\end{definition}

\section{Peripheral eigenvalue group structure}

This section proves the peripheral spectrum structure of~\cite[Theorem~6.6]{Wolf2012Quantum} for irreducible Schwarz maps:
the peripheral eigenvalues form a finite cyclic group.
The trace-preserving refinement that the order divides the bond
dimension is proved in
Theorem~\ref{thm:peripheral_cyclic_group}.

\begin{lemma}[Peripheral eigenvectors are invertible]%
    \label{lem:isUnit_peripheral_eigenvector}
    \lean{MPSTensor.isUnit_peripheral_eigenvector}
    \leanok
    \uses{thm:ks_peripheral}
    Let $E$ be an irreducible unital Kraus map on $\MN{D}$ with a
    positive-definite adjoint fixed point.
    If $X \ne 0$ satisfies $E(X) = \mu X$ with $|\mu| = 1$,
    then $X$ is invertible.
\end{lemma}

\begin{proof}
    \leanok
    The KS equality gives $E(X^\dagger X) = X^\dagger X$, so $X^\dagger X$
    is a nonzero PSD fixed point. Irreducibility upgrades it to positive
    definite, hence invertible, so $\det(X) \ne 0$.
\end{proof}

\begin{lemma}[Peripheral eigenvalues: product closure]%
    \label{lem:peripheral_mul_closure}
    \lean{MPSTensor.peripheralEigenvalues_mul_mem_of_irreducible_unital_of_adjoint_fixedPoint}
    \leanok
    \uses{lem:isUnit_peripheral_eigenvector,thm:ks_peripheral}
    Let $E$ be an irreducible unital Kraus map on $\MN{D}$ with a
    positive-definite adjoint fixed point.
    If $\mu, \nu$ are peripheral eigenvalues of $E$, then so is $\mu\nu$.
\end{lemma}

\begin{proof}
    \leanok
    Take eigenvectors $X, Y$ for $\mu, \nu$. The KS equality at $X$
    puts $X$ in the multiplicative domain, so
    $E(YX) = E(Y)E(X) = (\nu\mu)(YX)$. Since $X, Y$ are units
    (Lemma~\ref{lem:isUnit_peripheral_eigenvector}), $YX \ne 0$ and
    $|\mu\nu| = 1$.
\end{proof}

\begin{theorem}[Peripheral eigenvalues: cyclic structure]%
    \label{thm:peripheral_cyclic_structure}
    \lean{PeripheralSpectrum.peripheral_eigenvalues_cyclic_structure}
    \leanok
    \uses{thm:ks_peripheral}
    Let $E$ be an irreducible unital Schwarz map on $\MN{D}$ with a
    positive-definite adjoint fixed point (trace-preservation is \emph{not}
    required).
    Then there exist $m \ge 1$ and a primitive $m$-th root of unity
    $\gamma$ such that the peripheral eigenvalues of $E$ are exactly
    $\{1, \gamma, \gamma^2, \ldots, \gamma^{m-1}\}$.
\end{theorem}

\begin{proof}
    \leanok
    The peripheral eigenvalues form a finite subset of $\C$,
    contain $1$, and are closed under multiplication and inversion,
    so they form a finite subgroup of $\C^\times$.
    Any finite subgroup of $\C^\times$ is cyclic.
    The generator $\gamma$ has order $m = |\text{peripheral spectrum}|$.
\end{proof}

\begin{theorem}[Peripheral eigenvalues form a cyclic group]%
    \label{thm:peripheral_cyclic_group}
    \lean{PeripheralSpectrum.peripheral_eigenvalues_form_cyclic_group}
    \leanok
    \uses{thm:peripheral_cyclic_structure}
    Let $E$ be an irreducible unital trace-preserving Schwarz map on
    $\MN{D}$ with a positive-definite adjoint fixed point.
    Then there exist $m \ge 1$ and a primitive $m$-th root of unity
    $\gamma$ such that $m \mid D$ and the peripheral eigenvalues of $E$
    are exactly $\{1, \gamma, \gamma^2, \ldots, \gamma^{m-1}\}$.
\end{theorem}

\begin{proof}
    \leanok
    \uses{thm:peripheral_cyclic_structure,
        thm:cyclic_decomposition_irreducible_schwarz}
    Apply Theorem~\ref{thm:peripheral_cyclic_structure} for the cyclic
    structure. Divisibility $m \mid D$ follows from the cyclic
    decomposition: $m$ mutually orthogonal projections summing to $\Id$
    each have equal trace (by trace preservation), so
    $m \cdot \tr(P_0) = D$ with $\tr(P_0) \in \N$.
\end{proof}

\begin{theorem}[Period divides bond dimension]%
    \label{thm:channel_period_divides_dim}
    \lean{PeripheralSpectrum.channel_period_divides_dim}
    \leanok
    \uses{thm:cyclic_decomposition_irreducible_schwarz}
    Let $E$ be an irreducible unital trace-preserving Schwarz map on
    $\MN{D}$ with a positive-definite adjoint fixed point. Let $\gamma$
    be a primitive $m$-th root of unity such that the peripheral eigenvalues
    are exactly $\{\gamma^k : k \in \Fin{m}\}$.
    Then $m \mid D$.
\end{theorem}

\begin{proof}
    \leanok
    \uses{thm:cyclic_decomposition_irreducible_schwarz}
    The cyclic decomposition gives orthogonal projections
    $P_0,\ldots,P_{m-1}$ with $\sum_k P_k=\Id$ and
    $E(P_{k+1})=P_k$. Trace preservation gives
    $\tr(P_{k+1})=\tr(P_k)$ for every $k$, hence
    \[
        D=\tr(\Id)=\sum_{k=0}^{m-1}\tr(P_k)=m\,\tr(P_0).
    \]
    Since the trace of an orthogonal projection is an integer, $m\mid D$.
\end{proof}

\subsection{Cyclic decomposition of irreducible Schwarz maps}

\begin{definition}[Corner preservation]
    \lean{PreservesCorner}
    \leanok
    For an orthogonal projection $P$ and linear map $E$, the corner
    $P\MN{D}P$ is preserved if $E(PXP)=PE(PXP)P$ for all $X$.
\end{definition}

\begin{definition}[Corner subspace]
    \lean{cornerSubmodule}
    \leanok
    The corner subspace associated with $P$ is
    $\{PXP : X\in \MN{D}\}$.
\end{definition}

\begin{definition}[Restricted corner map]
    \lean{cornerRestriction}
    \leanok
    The restriction of $E$ to an invariant corner $P\MN{D}P$.
\end{definition}

\begin{definition}[Irreducible restriction on a corner]
    \lean{IsIrreducibleOnCorner}
    \leanok
    Irreducibility of the corner-restricted map.
\end{definition}

\begin{definition}[Rank of an orthogonal projection]
    \lean{cornerRank}
    \leanok
    The rank of an orthogonal projection $P \in \MN{D}$, i.e.\ the
    dimension of the corner algebra $P\MN{D}P$.
\end{definition}

\begin{lemma}[Compression isometry for an orthogonal projection]%
    \label{thm:compression_isometry_projection}
    \lean{cornerSubmoduleMatrixLinearEquiv}
    \leanok
    Let $P \in \MN{D}$ be an orthogonal projection of rank $n$. The corner
    algebra $P\MN{D}P$ is linearly isomorphic to the full matrix algebra
    $\MN{n}$. The isomorphism is built from the spectral diagonalisation
    of $P$ by conjugating the top-left $n\times n$ block by the unitary
    diagonalising $P$.
\end{lemma}

\begin{lemma}[Rank equals trace for an orthogonal projection]%
    \label{thm:corner_rank_eq_trace}
    \lean{cornerRank_eq_trace}
    \leanok
    The rank of an orthogonal projection $P$ equals its trace: $n = \tr P$.
\end{lemma}

\begin{lemma}[Scalar fixed points for irreducible unital maps]
    \label{lem:irreducible_unital_fixed_points_scalar}
    \lean{MPSTensor.fixed_eq_scalar_of_irreducible_unital}
    \leanok
    \uses{thm:psd_fp_unique_irred}
    If $E$ is irreducible and unital, then every fixed point of $E$ is a
    scalar multiple of $\Id$.
\end{lemma}

\begin{proof}
    \leanok
    \uses{thm:psd_fp_unique_irred}
    If $E(X)=X$, then the Hermitian matrices
    \[
        X+X^\dagger,
        \qquad
        i(X-X^\dagger)
    \]
    are also fixed by $E$. The positive-semidefinite fixed-point uniqueness
    theorem gives $X+X^\dagger=a\Id$ and
    $i(X-X^\dagger)=b\Id$, hence
    \[
        X=\frac{1}{2}(a-i b)\Id.
    \]
\end{proof}

\begin{lemma}[Peripheral unitary eigenvector]
    \label{lem:peripheral_unitary_eigenvector}
    \lean{MPSTensor.exists_peripheral_unitary_of_irreducible_schwarz}
    \leanok
    \uses{thm:ks_equality_of_peripheral_eigenvector_of_fixedPoint,
        thm:psd_fp_unique_irred}
    For an irreducible unital Schwarz map with faithful adjoint fixed point,
    each peripheral eigenvalue admits a unitary eigenvector.
\end{lemma}

\begin{lemma}[Powers on a peripheral-unitary orbit]
    \label{lem:peripheral_unitary_powers}
    \lean{MPSTensor.map_powers_of_peripheral_unitary}
    \leanok
    \uses{thm:ks_equality_of_peripheral_eigenvector_of_fixedPoint,
        thm:right_multiplicative_identity}
    If $E(U)=\mu U$ with $|\mu|=1$ and $U$ unitary, then
    $E(U^n)=\mu^n U^n$ for all $n$.
\end{lemma}

\begin{lemma}[Normalized peripheral unitary]
    \label{lem:normalized_peripheral_unitary}
    \lean{MPSTensor.exists_normalized_peripheral_unitary_of_irreducible_schwarz}
    \leanok
    \uses{lem:peripheral_unitary_eigenvector, lem:peripheral_unitary_powers,
        lem:irreducible_unital_fixed_points_scalar}
    Peripheral unitary eigenvectors can be normalized compatibly with
    the faithful fixed-point state.
\end{lemma}

\begin{lemma}[Cyclic projections from a peripheral unitary]
    \label{lem:cyclic_projections_from_peripheral_unitary}
    \lean{exists_cyclic_projections_of_peripheral_unitary}
    \leanok
    A primitive $m$-th peripheral root yields $m$ orthogonal projections
    summing to $\Id$ and cyclically permuted by the channel.
\end{lemma}

\begin{theorem}[Cyclic decomposition for irreducible Schwarz maps]
    \label{thm:cyclic_decomposition_irreducible_schwarz}
    \lean{MPSTensor.exists_cyclic_decomposition_of_irreducible_schwarz}
    \leanok
    \uses{lem:normalized_peripheral_unitary, lem:peripheral_unitary_powers,
        lem:cyclic_projections_from_peripheral_unitary}
    Under irreducibility and faithful fixed point hypotheses, there is a
    full cyclic projection decomposition realizing the peripheral period.
\end{theorem}

\begin{lemma}[Power maps preserve each cyclic sector]
    \lean{preserves_corner_pow_of_cyclic_decomp}
    \leanok
    Appropriate powers of the channel preserve each sector corner.
\end{lemma}

\begin{lemma}[Irreducibility of restricted sector dynamics]
    \lean{isIrreducible_restriction_of_cyclic_decomp}
    \leanok
    Each sector restriction is irreducible.
\end{lemma}

\begin{lemma}[Primitivity of restricted sector dynamics]
    \lean{isPrimitive_restriction_of_cyclic_decomp}
    \leanok
    Each sector restriction is primitive.
\end{lemma}

\begin{lemma}[Corner invariance at the period]
    \lean{preserves_corner_pow_orderOf_of_perm_decomp}
    \leanok
    The channel raised to the period preserves every cyclic corner.
\end{lemma}

\subsection{Group structure of the peripheral spectrum}

\begin{lemma}[Peripheral product closure]
    \lean{PeripheralSpectrum.peripheral_eigenvalues_closed_under_mul}
    \leanok
    \uses{lem:peripheral_mul_closure}
    The product of two peripheral eigenvalues is peripheral.
\end{lemma}

\begin{lemma}[Peripheral inverse closure]
    \lean{PeripheralSpectrum.peripheral_eigenvalues_closed_under_inv}
    \leanok
    \uses{thm:peripheral_cyclic_structure}
    The inverse of a peripheral eigenvalue is peripheral.
\end{lemma}

\begin{theorem}[Peripheral eigenvalue multiplicity one]
    \label{thm:peripheral_multiplicity_one}
    \lean{PeripheralSpectrum.peripheral_eigenvalue_multiplicity_one}
    \leanok
    \uses{lem:peripheral_unitary_eigenvector,
        thm:ks_equality_of_peripheral_eigenvector_of_fixedPoint,
        thm:right_multiplicative_identity,
        lem:irreducible_unital_fixed_points_scalar}
    Each peripheral eigenvalue has algebraic multiplicity one.
\end{theorem}
\begin{proof}\leanok
    Given a peripheral unitary eigenvector $U$ for $\gamma$, the
    multiplicative-domain identity gives $E(X U^\dagger) = X U^\dagger$
    for any $\gamma$-eigenvector $X$. By the scalar fixed-point lemma
    for irreducible unital maps, $X U^\dagger = c \cdot \Id$, hence
    $X = c \cdot U$ and the eigenspace is one-dimensional.
\end{proof}

\section{Primitive overlap convergence (complementary transfer-map gap form)}

The theorems in this section are stated directly in terms of the
complementary map $E - P$. This is the analytic form needed for the later
overlap argument, corresponding to the primitive-case specialization
of~\cite[Theorem~6.7]{Wolf2012Quantum} where $T_\phi$ is a rank-one
projection.
For primitive normalized transfer maps, Chapter~\ref{ch:channels} supplies
this complementary transfer-map input under explicit hypotheses.

\begin{theorem}[Complementary transfer-map gap implies trace convergence to $1$]%
    \label{thm:trace_pow_one}
    \lean{MPSTensor.linearMap_trace_pow_tendsto_one_of_spectralRadius_compl_lt_one}
    \leanok
    \uses{def:fixed_point_proj, def:trace_preserving}
    Let $E$ be trace-preserving with fixed point $\rho$
    ($\tr(\rho) \neq 0$) and
    let $P$ be the fixed-point projection
    (Definition~\ref{def:fixed_point_proj}).
    If $\rho_{\spec}(E - P) < 1$, then
    $\Tr(E^n) \to 1$ as $n \to \infty$.
\end{theorem}

\begin{proof}
    \leanok
    \uses{thm:pow_decomp}
    By the power decomposition
    (Theorem~\ref{thm:pow_decomp}),
    $E^n = P + (E - P)^n$ for $n \ge 1$.
    Since $\rho_{\spec}(E - P) < 1$, the Gelfand formula gives
    $(E - P)^n \to 0$, so $\Tr(E^n) \to \Tr(P)$.
    The operator trace of $P$ equals~$1$: since
    $P(X) = \frac{\tr(X)}{\tr(\rho)} \rho$,
    we have $\Tr(P) = \frac{\tr(\rho)}{\tr(\rho)} = 1$.
\end{proof}

\begin{theorem}[Self-overlap convergence from a complementary transfer-map gap]%
    \label{thm:primitive_overlap}
    \lean{MPSTensor.mpvOverlap_tendsto_one_of_transfer_spectralRadius_compl_lt_one}
    \leanok
    \uses{def:transfer_map, def:fixed_point_proj, def:mpv_overlap}
    Let $A$ be a normalized MPS tensor with a nonzero PSD fixed
    point $\rho$ of the transfer map.
    If $\rho_{\spec}(\E_A - P) < 1$ (where $P$ is the fixed-point
    projection), then $O_{AA}(N) \to 1$ as $N \to \infty$.
\end{theorem}

\begin{proof}
    \leanok
    \uses{thm:trace_pow_one, thm:overlap_trace}
    By Lemma~\ref{thm:overlap_trace} (with $A = B$),
    $O_{AA}(N) = \Tr(\E_A^N)$.
    By Theorem~\ref{thm:trace_pow_one},
    $\Tr(\E_A^N) \to 1$.
\end{proof}

\begin{lemma}[Norm of self-overlap convergence]\label{lem:norm_selfOverlap_tendsto}
    \lean{MPSTensor.tendsto_norm_selfOverlap_one}\leanok
    \uses{def:mpv_overlap}
    If $A$ is an MPS tensor and $O_{AA}(N) \to 1$ in $\C$ as
    $N \to \infty$, then $\|O_{AA}(N)\| \to 1$ as
    $N \to \infty$.
\end{lemma}

\begin{proof}\leanok
    Follows from the continuity of the norm:
    $O_{AA}(N) \to 1$ implies
    $\|O_{AA}(N)\| \to \|1\| = 1$.
\end{proof}
